% GENERATED FILE. Edit canonical lessons, then run npm run generate:books.
% canonical-source-hash: fnv1a64:d724db6bfb946e8c
% canonical-lessons: GE-C183-leiden, GE-C183-messen, GE-C183-mischen, GE-C183-naehen, GE-C183-niesen

\chapter{To suffer, To measure, To mix, To sew, To sneeze}
\label{ch:ge-to-suffer-to-measure-to-mix-to-sew-to-sneeze}
\hlchaptermodality{183}

\begin{chapteropening}
\textbf{By the end of this chapter:} \emph{I can say to suffer, to measure, to mix, to sew and to sneeze.}

Complete the last lesson of chapter 183: I can say to suffer, to measure, to mix, to sew and to sneeze.
\end{chapteropening}

\section[leiden]{leiden — to suffer}

\label{lesson:GE-C183-leiden}



\begin{quote}
Before the new one: say the German for to land, then the German for to smile.
\end{quote}

\subsection*{You'll want to know: leiden}
\textbf{leiden} — \textquotedblleft{}to suffer\textquotedblright{}.

\begin{cousinweb}
\textbf{leiden} $\leftarrow$ Middle High German \emph{līden} $\leftarrow$ Old High German \emph{līdan}, \textquotedblleft{}to go, to undergo\textquotedblright{} (EtymDB 2.1) — inherited. To suffer is to go through something. Dutch \emph{lijden} is the same verb.
\end{cousinweb}

\begin{grammarlens}[title={What you've built}]
One more word for this chapter.
\end{grammarlens}

\subsection*{Your turn}
\begin{itemize}
\raggedright
  \item \emph{Say it:} \emph{leiden}
  \item \emph{Say it:} \emph{leiden}, once more
  \item \emph{Say it:} say \emph{leiden}
  \item \emph{Recall:} say the German for to land, then the German for to smile, then say \emph{leiden} again
\end{itemize}

\begin{tcolorbox}[breakable,colback=teal!4,colframe=teal!35!black,title={Before you move on}]
What does leiden mean? (\textquotedblleft{}To suffer\textquotedblright{}.) Say it once more.
\end{tcolorbox}
\section[messen]{messen — to measure}

\label{lesson:GE-C183-messen}



\begin{quote}
Before the new one: say the German for to smile, then the German for to suffer.
\end{quote}

\subsection*{You'll want to know: messen}
\textbf{messen} — \textquotedblleft{}to measure\textquotedblright{}.

\begin{cousinweb}
\textbf{messen} $\leftarrow$ Old High German \emph{mezzen} $\leftarrow$ Proto-Germanic \emph{*metaną}, \textquotedblleft{}to measure\textquotedblright{} — inherited, the same verb as English \emph{mete}, as in \textquotedblleft{}mete out\textquotedblright{}. (EtymDB 2.1)
\end{cousinweb}

\begin{grammarlens}[title={What you've built}]
One more word for this chapter.
\end{grammarlens}

\subsection*{Your turn}
\begin{itemize}
\raggedright
  \item \emph{Say it:} \emph{messen}
  \item \emph{Say it:} \emph{messen}, once more
  \item \emph{Say it:} say \emph{messen}
  \item \emph{Recall:} say the German for to smile, then the German for to suffer, then say \emph{messen} again
\end{itemize}

\begin{tcolorbox}[breakable,colback=teal!4,colframe=teal!35!black,title={Before you move on}]
What does messen mean? (\textquotedblleft{}To measure\textquotedblright{}.) Say it once more.
\end{tcolorbox}
\section[mischen]{mischen — to mix}

\label{lesson:GE-C183-mischen}



\begin{quote}
Before the new one: say the German for to suffer, then the German for to measure.
\end{quote}

\subsection*{You'll want to know: mischen}
\textbf{mischen} — \textquotedblleft{}to mix\textquotedblright{}.

\begin{cousinweb}
\textbf{mischen} $\leftarrow$ Old High German \emph{miskan}, from Latin \emph{miscēre}, \textquotedblleft{}to mix\textquotedblright{} — an early loan. EtymDB 2.1 has the Latin verb.
\end{cousinweb}

\begin{grammarlens}[title={What you've built}]
One more word for this chapter.
\end{grammarlens}

\subsection*{Your turn}
\begin{itemize}
\raggedright
  \item \emph{Say it:} \emph{mischen}
  \item \emph{Say it:} \emph{mischen}, once more
  \item \emph{Say it:} say \emph{mischen}
  \item \emph{Recall:} say the German for to suffer, then the German for to measure, then say \emph{mischen} again
\end{itemize}

\begin{tcolorbox}[breakable,colback=teal!4,colframe=teal!35!black,title={Before you move on}]
What does mischen mean? (\textquotedblleft{}To mix\textquotedblright{}.) Say it once more.
\end{tcolorbox}
\section[nähen]{nähen — to sew}

\label{lesson:GE-C183-naehen}



\begin{quote}
Before the new one: say the German for to measure, then the German for to mix.
\end{quote}

\subsection*{You'll want to know: nähen}
\textbf{nähen} — \textquotedblleft{}to sew\textquotedblright{}.

\begin{cousinweb}
\textbf{nähen} $\leftarrow$ Old High German \emph{nājen} $\leftarrow$ Proto-Germanic \emph{*nēaną}, \textquotedblleft{}to sew, to stitch\textquotedblright{} — inherited. (EtymDB 2.1) The German word for \textquotedblleft{}needle\textquotedblright{} comes from the same root.
\end{cousinweb}

\begin{grammarlens}[title={What you've built}]
One more word for this chapter.
\end{grammarlens}

\subsection*{Your turn}
\begin{itemize}
\raggedright
  \item \emph{Say it:} \emph{nähen}
  \item \emph{Say it:} \emph{nähen}, once more
  \item \emph{Say it:} say \emph{nähen}
  \item \emph{Recall:} say the German for to measure, then the German for to mix, then say \emph{nähen} again
\end{itemize}

\begin{tcolorbox}[breakable,colback=teal!4,colframe=teal!35!black,title={Before you move on}]
What does nähen mean? (\textquotedblleft{}To sew\textquotedblright{}.) Say it once more.
\end{tcolorbox}
\section[niesen]{niesen — to sneeze}

\label{lesson:GE-C183-niesen}



\begin{quote}
Before the new one: say the German for to mix, then the German for to sew.
\end{quote}

\subsection*{You'll want to know: niesen}
\textbf{niesen} — \textquotedblleft{}to sneeze\textquotedblright{}.

\begin{cousinweb}
\textbf{niesen} $\leftarrow$ Old High German \emph{niosan} $\leftarrow$ Proto-Germanic \emph{*hneusaną}, \textquotedblleft{}to sneeze\textquotedblright{} — inherited. (EtymDB 2.1) Dutch \emph{niezen} is the same verb.
\end{cousinweb}

\begin{grammarlens}[title={What you've built}]
One more word for this chapter.
\end{grammarlens}

\subsection*{Your turn}
\begin{itemize}
\raggedright
  \item \emph{Say it:} \emph{niesen}
  \item \emph{Say it:} \emph{niesen}, once more
  \item \emph{Say it:} say \emph{niesen}
  \item \emph{Recall:} say the German for to mix, then the German for to sew, then say \emph{niesen} again
\end{itemize}

\begin{tcolorbox}[breakable,colback=teal!4,colframe=teal!35!black,title={Before you move on}]
What does niesen mean? (\textquotedblleft{}To sneeze\textquotedblright{}.) Say it once more.
\end{tcolorbox}
